\epanalipsi{Α3}

\tmima{Λεξικό του κεφαλαίου}

\renewcommand{\arraystretch}{1.5}
\begin{longtable}{@{}>{\raggedright\arraybackslash}p{0.28\linewidth} >{\raggedright\arraybackslash}p{0.62\linewidth}@{}}
\toprule
\textbf{Λέξη ή τύπος} & \textbf{Με μία πρόταση}\\
\midrule
\endhead
ρητοί & Οι αριθμοί που γράφονται ως κλάσμα $\frac{\alpha}{\beta}$ με $\alpha,\beta$ ακέραιους και $\beta\neq0$.\\
ισοδύναμα κλάσματα & Ίδια αξία: $\frac{\alpha}{\beta}=\frac{\gamma}{\delta}$ όταν $\alpha\delta=\beta\gamma$.\\
ανάγωγο & Κλάσμα που δεν απλοποιείται: $\MKD$ των όρων ίσος με $1$.\\
$|\alpha|$ & Η απόσταση του $\alpha$ από το $0$: ποτέ αρνητική.\\
αντίθετοι & Ετερόσημοι με ίδια απόλυτη τιμή: $\alpha$ και $-\alpha$, με άθροισμα $0$.\\
σύγκριση & Μεγαλύτερος είναι ο δεξιότερος στον άξονα.\\
$\alpha-\beta$ & $=\alpha+(-\beta)$: προσθέτω τον αντίθετο.\\
γινόμενο, πηλίκο & Ομόσημοι: θετικό. Ετερόσημοι: αρνητικό. Ποτέ διαίρεση με το $0$.\\
αντίστροφος & Ο $\frac{\beta}{\alpha}$ για τον $\frac{\alpha}{\beta}$: γινόμενο $1$. Το $0$ δεν έχει.\\
$\alpha^{\nu}$ & $\nu$ παράγοντες ίσοι με $\alpha$. Ίδιο πρόσημο με το γινόμενο.\\
$\alpha\cdot10^{\nu}$ & Τυποποιημένη μορφή: $1\le\alpha<10$ και $\nu$ μη μηδενικός φυσικός.\\
\bottomrule
\end{longtable}

\tmima{Ασκήσεις επανάληψης}

\ex Για καθέναν από τους $-\frac{6}{3}$, $\frac{4}{5}$, $0$, $-2{,}5$, $\frac{7}{7}$ πες αν είναι φυσικός, ακέραιος, ρητός.

\ex Απλοποίησε το $\frac{36}{48}$ μέχρι να γίνει ανάγωγο και γράψε ένα κλάσμα ισοδύναμο με το $\frac{5}{6}$ που έχει παρονομαστή $42$.

\ex Γράψε το $\frac{7}{8}$ ως δεκαδικό και τον $0{,}36$ ως ανάγωγο κλάσμα.

\ex Βρες: α) $|-7{,}5|$ \quad β) $\left|\frac{3}{4}\right|$ \quad γ) τον αντίθετο του $-\frac{2}{3}$ \quad δ) όλους τους $\alpha$ με $|\alpha|=4$.

\ex Γράψε σε αύξουσα σειρά τους αριθμούς $-2{,}5$, $0$, $-\frac{1}{2}$, $\frac{3}{2}$, $-3$, $1{,}2$.

\ex Βρες δύο ρητούς ανάμεσα στους $\frac{2}{7}$ και $\frac{3}{7}$.

\ex Υπολόγισε: α) $-\frac{3}{4}+\frac{5}{6}$ \quad β) $2-\frac{7}{5}$ \quad γ) $-1{,}5+0{,}75$ \quad δ) $\frac{1}{2}-\left(-\frac{1}{3}\right)$.

\ex Υπολόγισε την παράσταση $-\frac{1}{2}+\frac{2}{3}-\frac{1}{6}+\frac{1}{4}$.

\ex Υπολόγισε: α) $\left(-\frac{3}{4}\right)\cdot\frac{8}{9}$ \quad β) $(-0{,}5)\cdot(-1{,}2)$ \quad γ) $12\cdot\left(\frac{1}{3}+\frac{1}{4}\right)$.

\ex Υπολόγισε: α) $\left(-\frac{3}{5}\right):\frac{9}{10}$ \quad β) $6:\left(-\frac{3}{4}\right)$ \quad γ) $7{,}2:0{,}09$.

\ex Μετάτρεψε σε απλό κλάσμα το $\dfrac{\frac{2}{3}+\frac{1}{6}}{\frac{5}{4}}$.

\ex Υπολόγισε: α) $\left(-\frac{2}{3}\right)^3$ \quad β) $-3^2$ \quad γ) $(-0{,}1)^2$ \quad δ) $\frac{1}{2}+\left(-\frac{1}{2}\right)^2$.

\ex Γράψε σε τυποποιημένη μορφή τον $45\,000\,000$, σε πλήρη μορφή τον $3{,}02\cdot10^{5}$ και σύγκρινε τους $8\cdot10^{6}$ και $9{,}5\cdot10^{5}$.

\ex \textbf{Ο αγώνας.} Ένα σχολείο έχει $240$ μαθητές. Τα $\frac{3}{8}$ παίζουν μπάσκετ και το $\frac{1}{3}$ αυτών είναι κορίτσια. Πόσα κορίτσια παίζουν μπάσκετ;

\ex \textbf{Σωστό ή λάθος (πράξεις).} α) $-\frac{3}{4}<-\frac{2}{3}$. \quad β) $\frac{1}{3}+\frac{1}{3}=\frac{2}{6}$. \quad γ) $(-2)\cdot(-3)\cdot(-4)=-24$. \quad δ) $5:\frac{1}{5}=1$. \quad ε) Το $0$ έχει αντίστροφο. \quad στ) Ο αντίθετος του $0$ είναι το $0$.

\ex \textbf{Σωστό ή λάθος (δυνάμεις, τυποποιημένη μορφή).} α) $(-2)^3=-6$. \quad β) $-2^2=-4$. \quad γ) $\left(\frac{2}{3}\right)^2=\frac{4}{3}$. \quad δ) $5\cdot10^{4}=50\,000$. \quad ε) Το $12\cdot10^{3}$ είναι σε τυποποιημένη μορφή. \quad στ) Τα $\frac{3}{4}$ και $\frac{9}{12}$ είναι ισοδύναμα.

\Needspace{24\baselineskip}\tmima{Ο χάρτης της επανάληψης}

Μετά τις λύσεις, δες σε ποιες ασκήσεις δυσκολεύτηκες. Ο πίνακας σου λέει ποια παράγραφο να ξαναδιαβάσεις.

\renewcommand{\arraystretch}{1.5}
\begin{tabular*}{\linewidth}{@{\extracolsep{\fill}} >{\raggedright\arraybackslash}p{0.38\linewidth} >{\raggedright\arraybackslash}p{0.52\linewidth}}
\toprule
\textbf{Αν δυσκολεύτηκες στις} & \textbf{ξαναδιάβασε την παράγραφο}\\
\midrule
ΕΑ3.1, ΕΑ3.2, ΕΑ3.3 & Α3.1 (ρητοί, ισοδύναμα, δεκαδικοί)\\
ΕΑ3.4, ΕΑ3.5, ΕΑ3.6 & Α3.2 (απόλυτη τιμή, διάταξη, ενδιάμεσοι ρητοί)\\
ΕΑ3.7, ΕΑ3.8 & Α3.3 (πρόσθεση και αφαίρεση)\\
ΕΑ3.9, ΕΑ3.14 & Α3.4 (πολλαπλασιασμός, κλάσμα ποσότητας)\\
ΕΑ3.10, ΕΑ3.11 & Α3.5 (διαίρεση, σύνθετα κλάσματα)\\
ΕΑ3.12 & Α3.6 (δυνάμεις)\\
ΕΑ3.13 & Α3.7 (τυποποιημένη μορφή)\\
ΕΑ3.15, ΕΑ3.16 & Α3.1 έως Α3.7\\
\bottomrule
\end{tabular*}
